\section{Modular lifting in the orders congruent to four modulo eight}
\label{sec:four-mod-eight-lifting}

The preceding lifting theorem used $N=4(q+1)$ with odd $q$, so its
orders were divisible by eight. The complementary choice $N=4q$
gives an especially simple result throughout the other admissible
congruence class. Its only spurious modular extensions, once five
orthogonal rows are present, are the rows already installed.
This is a counting and lifting statement; it does not guarantee the
surplus modular solutions needed for a new row. No novelty claim is
made pending a targeted prior-art check.

Let $P$ have $r$ orthogonal sign rows $p_i$ of length $N$. Count sign
vectors up to negation, and write $M_q(P)$ for the number satisfying
$Pv\equiv0\pmod q$ and $E(P)$ for the number satisfying $Pv=0$.

\begin{theorem}[The row correction at orders $4\bmod8$]
Let $q\ge3$ be any odd integer, $N=4q$, and $r\ge5$. Then
\[
 \boxed{E(P)=M_q(P)-r.}
\]
Every modular sign extension which is not exact is the sign or
negative of exactly one existing row. Primality of $q$ is not required.
\end{theorem}
\begin{proof}
For a modular sign vector $v$ put $d_i=p_i\cdot v$.
Orthogonal rows differ in $N/2=2q$ positions, an even number, so
all $d_i$ have the same residue modulo four. Every $d_i$ is even
and divisible by odd $q$.

If that common residue is two, each $d_i$ must be $2q$ or $-2q$,
since the next possible magnitude $6q$ exceeds $N$. But Bessel gives
\[
 4q^2r=\sum_i d_i^2\le N\|v\|^2=N^2=16q^2,
\]
contradicting $r\ge5$. Hence all $d_i$ are zero modulo four and
therefore belong to $\{0,N,-N\}$. A nonzero value means
$v=\pm p_i$, by equality in the sign-vector inner product bound.
All other inner products then vanish by row orthogonality.
Conversely every existing row is a modular sign vector because its
nonzero self-inner-product is $N=4q$. The $r$ row classes are distinct,
proving the count formula.
\end{proof}

\subsection{The sharp four-row boundary and closed quadruples}
The five-row threshold has an exact geometric explanation. Call four
orthogonal sign rows a \emph{closed quadruple} if their entrywise
product is constant. This is standard terminology in Hadamard
switching; see Orrick, \emph{Switching Operations for Hadamard
Matrices}, Section 3.1 and Proposition 3.1,
\url{https://arxiv.org/abs/math/0507515}.
In particular the prohibition of such a quadruple inside a full
Hadamard matrix of order $4\bmod8$ is classical.

\begin{proposition}[Exact boundary correction]
For four orthogonal sign rows of length $N=4q$, $q\ge3$ odd,
\[
 E(P)=M_q(P)-4-4\,\mathbf1_{\{P\text{ is closed}\}}.
\]
When the quadruple is closed, the four extra modular classes are
represented by half-sums
\[
 v=\frac12\sum_{i=1}^4\sigma_i p_i,
 \qquad\sigma_i\in\{\pm1\},\qquad
 \prod_i\sigma_i=-c,
\]
where $c$ is the constant entrywise row product. Negating all four
$\sigma_i$ gives the same antipodal class.
\end{proposition}
\begin{proof}
The common-zero residue case in the preceding proof still produces
only exact extensions and existing rows. In the common-two case,
all four inner products are $\sigma_iN/2$, and Bessel is saturated.
Thus $v$ has exactly the stated half-sum representation, with no
orthogonal-complement component. Four signs sum to $+2$ or $-2$
exactly when their product is $-1$. Consequently every coordinate
of this half-sum is a sign if and only if
$\prod_i p_{ij}=-\prod_i\sigma_i$ for every coordinate $j$.
This is exactly the closed condition with the displayed restriction.
Eight choices of $\sigma$ modulo simultaneous negation give four
classes; independence of the four rows makes their half-sums distinct.
\end{proof}

For example $H_4\otimes\mathbf1_q$ is a closed quadruple with no
exact extension, yet has eight modular classes: four installed rows
and four half-sums. Thus subtracting only the original rows would
give a false positive at this sharp boundary. The no-extension claim
follows directly from its four length-$q$ block sums: orthogonality
forces all sums zero, impossible for odd block size.

\subsection{Finite algebra and the remaining positivity problem}
For any integer modulus $q$, use the integral group algebra of
$(\mathbb Z/q\mathbb Z)^r$ and let $c_j$ be column $j$ of $P$.
Then
\[
 C_q(P)=[X^0]\prod_{j=1}^N(X^{c_j}+X^{-c_j})=2M_q(P).
\]
Under the five-row hypothesis,
\[
 \boxed{E(P)=\tfrac12C_q(P)-r.}
\]
Hence $C_q(P)>2r$ is an exact sufficient condition for a new sign row,
and $C_q(P)=2r$ certifies terminality. The group-algebra formula
works for composite odd $q$ as well as primes. For prime $q$ it can
also be computed by finite-field Gaussian elimination or a Fourier
character sum. When composite moduli are used, ordinary field row
reduction must be replaced by valid arithmetic over the residue ring
or a suitable Chinese-remainder procedure.

For prime $q$, the row code $C\subseteq\mathbb F_q^N$ generated by
$P$ is self-orthogonal, because $PP^{\mathsf T}=NI_r\equiv0\pmod q$.
If its rank is $s$, then $s\le N/2$. An information-coordinate search
therefore uses $2^{N-s-1}$ sign assignments; large real row count
does not imply equally large modular rank. The existing rows already
give $M_q(P)\ge r$. Proving the strict surplus needed for extension
requires additional code structure, rather than rank alone.

Exact examples from explicit Hadamard matrices illustrate the formula:
\[
\begin{array}{c|r|r|r|r}
N&r&\operatorname{rank}_{q}P&M_q(P)&E(P)\\\hline
12&5&5&18&13\\
12&6&6&12&6\\
12&10&6&12&2\\
20&10&10&20&10\\
20&18&10&20&2\\
28&14&14&28&14\\
28&26&14&28&2
\end{array}
\]
Here $q=N/4$ and the matrices are constructed explicitly in the
standard-library checker. An independent integer group-algebra
calculation gives $C_3(P)=36$ for the five-row order-twelve example,
so $36/2-5=13$ exactly. Its thirteen candidate classes need not all
be mutually orthogonal: counting single-row extensions alone does
not automatically supply a complete frame.

Together, the $N=4(q+1)$ theorem for odd $q$ and the present
$N=4q$ theorem treat both admissible residue classes modulo eight,
with different explicit corrections. Neither asserts the requisite
modular positivity uniformly, so neither proves the full Hadamard
conjecture. Their concrete value is that they replace an approximate
modular extension problem by an exact real one after accounting for
fully classified exceptional configurations.
